\section{Hermite Polynomials}

Consider $q^{(2)}_n$ when $\wf^{(G)}(x) \deff (\sqrt{2\pi})^{-1} \ee^{-\frac{1}{2} x^2}$. We make standard the notation
\begin{equation}
	\He_n(x) \deff q^{(2,\wf^{(G)})}_n.
\end{equation}
This means that \textit{Hermite Polynomials} are such that
\begin{equation}
	\int_{-\infty}^{\infty} \He_n(x) \He_m(x) \wf^{(G)}(x) \dd x = 
	\begin{cases}
		0, \quad \text{if} \quad n \neq m \\
		n!, \quad \text{if} \quad n = m.
	\end{cases}
\end{equation}
There are some other representation for the polynomial - one of the most common is their derivative formula
\begin{equation}
	\He_n(x) = (-1)^{(n)}\ee^{\frac{1}{2} x^2} \frac{\dd}{\dd x^n} \ee^{-\frac{1}{2} x^2}
\end{equation}
the recursion formula 
\begin{equation}
	\He_{n+1}(x) = x \He_n(x) - n \He_{n-1}(x).
\end{equation}
or even the differential equation
\begin{equation}
	\He''_n(x) - 2 x \He'_n(x) = -2 n \He_n(x)
	\label{Equation: Differential equation}
\end{equation}
One can also write the explicit representation (or \textit{Standart}),
\begin{equation}
	\He_{n}(x) = n! \sum_{l=0}^{\lfloor \frac{n}{2} \rfloor} (-1)^{l} \frac{(x)^{n-2l}}{2^l (n-2l)! l!}.
	\label{Equation: Standart Hermite Notation}
\end{equation}

This is all good, but can be a little inconvenient when doing computation. A nice representation is presented by \cite{nica2025gaussianintegralformulahermite}. The Gaussian integral formula for the Hermite Polynomials. This development will be based on this reference.
\begin{definition}
	The Gaussian integral Formula, or equivalently the Gaussian Expectation formula, for the Hermite polynomials $\He_n$ is given by
	\begin{equation}
		\He_n(x) = \E_{Z \sim \mathcal{N}(0,1)}[(x+\ii Z)^n] = \int_{-\infty}^{\infty} (x+\ii z)^n \ee^{-\frac{1}{2} z^2} \dd z
	\end{equation}
	\label{Def: Gaussian integral-expectation formula}
\end{definition}
This deeply simplifies some computation. For example, note that someone can easily compute the identities
\begin{align}
	\He_n(0) & = \E[(0 + \ii Z)^n] = \begin{cases}
		(-1)^{n/2}(n-1)!!, \quad \text{for $n$ even} \\
		0, \quad \text{for $n$ odd}
	\end{cases}, \\
	\He_n(x+y) & = \E[(x+y+\ii Z)^n] = \E\left[\sum_{k=0}^{n} \binom{n}{k} y^k (x+\ii Z)^{n-k}\right] = \sum_{k=0}^{n} \binom{n}{k} y^k \He_{n-k}(x).
\end{align}
And with that one can get a explicit expression for the polynomial setting $x=0$ in the second expression and using the first one to calculate the values.

\subsection{Arbitrary Variance}

We can use this formula to work with Hermite Polynomials of arbitrary variance. Assume one adds a temporal dependence parameter $t$. Denote this polynomial $\Het_n(x;t)$ defined as
\begin{definition}
	\begin{equation}
		\Het_n(x;t) \deff \E_{B_t \sim \mathcal{N}(0,t)}\left[(x+\ii B_t)^n\right] = \E_{Z \sim \mathcal{N}(0,1)}\left[(x+\ii \sqrt{t} Z)^n\right]
	\end{equation}
\end{definition}
The notation has the tilde because historically there is a preference for the derivative formula instead of the integral formula. The notation $\Her_n$ is reserved for the dual of this polynomial, defined as below
\begin{definition}
	\begin{equation}
		\Her_n(x;t) \deff \E_{Z \sim \mathcal{N}(0,\frac{1}{t})}\left[(\frac{x}{t}+\ii Z)^n\right] = \E_{Z \sim \mathcal{N}(0,1)}\left[(\frac{x}{t}+\ii \frac{Z}{\sqrt{t}})^n\right].
	\end{equation}
	Which is dual in the sense that
	\begin{equation}
		\int_{-\infty}^{\infty} \Het_n(x;t) \Her_m(x;t) (\sqrt{2\pi t})^{-1} \ee^{-\frac{1}{2} \frac{x^2}{t}} \dd x = 
		\begin{cases}
			0, \quad \text{if} \quad n \neq m \\
			n!, \quad \text{if} \quad n = m.
		\end{cases}
	\end{equation}
\end{definition}
Finally, note that
\begin{equation}
	\Het_n(x;t) = t^{n/2} \Het\left(\frac{x}{\sqrt{t}};1\right) = t^{n/2} \He\left(\frac{x}{\sqrt{t}}\right), \quad \Her_n(x;t) = t^{-n/2} \Het\left(\frac{x}{\sqrt{t}};1\right) = t^{-n/2} \He\left(\frac{x}{\sqrt{t}}\right).
\end{equation}

\subsection{Combinatorial Arguments} 

Our attention will be on the Asymptotic of the Hermite Polynomials, to get there, we will work our way with some combinatorial arguments. We start with the \textit{Wick Formula/Isserlis' Theorem} for Gaussian expectations.
\begin{thm}
	(Wick Formula/Isserlis' Theorem) If $\mmany{Z}{n}$ are a collection of mean zero, jointly Gaussian random variables, then the expectation of the product $Z_1 \cdots Z_n$ can be written in terms of the individual correlations $\E[Z_a Z_b]$ in the following way
	\begin{equation}
		\E[Z_1 \cdots Z_n] = \sum_{\pi \in \Prt(\{1,\cdots,n\})} \prod_{\{a,b\} \in \pi} \E[Z_a Z_b]
	\end{equation}
	where $\Prt(\{1,\cdots,n\})$ are the set of pairings (or pair partitions) on the set $S = \{1, \cdots, n\}$. That is, $\Prt(S)$ denotes the way of dividing $S$ into disjoint pairs.
	\label{Theorem: Wick-Isserlis}
\end{thm}

If one applies this result for example, for $\E[Z^n]$, the expectation of a $n$-repeated Gaussian Variable, this result gives the moments of the normal distribution
\begin{equation}
	\E[Z^n] = |\Prt(\{1,\cdots,n\})| \cdot \E[Z^2]^{n/2} = 
	\begin{cases}
		(n-1)!! \cdot \Var[Z]^{n/2}, \quad \text{if $n$ is even}, \\
		0, \quad \text{if $n$ is odd}.
	\end{cases}
	\label{Equation: Normal Moments}
\end{equation}

If we want to apply this combinatorial argument for the polynomials $\Het_n(x;t) = \E_{Z \sim \mathcal{N}(0,1)}[(x+\ii \sqrt{t}Z)^n]$, we must expand the product in one of the terms on each of the brackets - which yields a sum over $2^n$ monomials. Then, taking the expectation over $Z$ amounts to adding up the pairings for each term as in \ref{Theorem: Wick-Isserlis}. The number of pairings depends on the number of times $m$ that $\ii \sqrt{t} Z$ was chosen: there are $(m-1)!!$ such pairings if $m$ is even and $0$ pairings if $m$ is odd. Therefore, the only contributions come from the terms with even order.

When exactly $m$ factors $\ii \sqrt{t} Z$ are chosen, there must be $n-m$ copies of $x$ chosen, and so the resulting contribution of the term is $(\ii \sqrt{t})^m x^{n-m}$. Since $m$ is even we will actually have $(-t)^{m/2} x^{n-m}$. We can count $\binom{n}{m}(m-1)!!$ such terms - $m$ choices of the term and the subsequent $(m-1)!!$ ways of pairing it up. This can be formalized as
\begin{align}
	\Het_n(x;t) & = \sum_{S \subset \{1,\cdots,n\}, \pi \in \Prt(S)} (\ii \sqrt{t})^{|\pi|} x^{n-|S|} \\
	& = \sum_{\{m=0 | m \ \text{is even}\}} \binom{n}{m} (m-1)!! (-t)^{m/2} x^{n-m} \\
	& = \sum_{k=0}^{\floor{\frac{n}{2}}} \binom{n}{2k} (2k-1)!! (-t)^{k} x^{n-2k} \\
	& = x^n - \binom{n}{2} 1!! \cdot t x^{n-2} + \binom{n}{4} 3!! \cdot t^2 x^{n-4} - \cdots
\end{align} 

This property can also be used to prove the orthogonality and recursion of the polynomials - for the proof, check \cite{nica2025gaussianintegralformulahermite}.

\subsection{Asymptotic - $x = \Boh(n)$}
\label{Subsection: Hermite x=n}

We can study the asymptotic of this polynomial for $x = \Boh(n)$ only using the combinatorial argument given in this chapter.

\begin{prop}
	Let $c \in \R$ be fixed, and set $x = c n$. We have the following asymptotic formula as $n \rightarrow \infty$:
	\begin{equation}
		\He_n(cn) = (cn)^n \ee^{-\frac{1}{2c^2}} (1 + \boh(1))
	\end{equation}
	that is,
	\begin{equation}
		\He_n(cn) \approx x^n \ee^{-\frac{n^2}{2x^2}}
	\end{equation}
	for $x = \Boh(n)$.
	\label{Prop: Hermite Asymptotic o(n)}
\end{prop}
\begin{proof}
	Using the integral formula \ref{Def: Gaussian integral-expectation formula} we can write
	\begin{equation}
		\He_n(cn) = \E[(cn + \ii Z)^n] = (cn)^n \E\left[\left(1 + \ii\frac{Z}{cn} \right)^n \right].
	\end{equation}
	Now, note that 
	\begin{equation}
		\lim_{n\rightarrow \infty} \left(1 + \ii\frac{Z}{cn} \right)^n = \ee^{\ii \frac{Z}{c}}
	\end{equation}
	almost surely in $Z$, and in light of the expectation $\E[\ee^{\beta Z}] = \ee^{\frac{1}{2} \beta^2}$we find that 
	\begin{equation}
		\E\left[\left(1 + \ii\frac{Z}{cn} \right)^n \right] = \ee^{-\frac{1}{2c^2}}(1+ \boh(1))
	\end{equation}
	from where the results follows.
\end{proof}

Other regimes for the asymptotic will have to use some more sophisticated techniques.

\subsection{Integral Representation}
\label{Subsection: Integral Hermite}


Follow \cite[Proposition 2.1]{Zhipeng-Liu-ClassNotes}. A standard method of asymptotic of orthogonal polynomials is the so called \textit{Deift-Zhou} steepest descent method, which considers to asymptotically solve a Riemann-Hilbert problem associated to the orthogonal polynomials. However, the asymptotic may me simpler to compute if the polynomial has an integral representation by direct inspection. The Hermite polynomials have such representation.

\begin{prop}
	The Hermite polynomial $\He_{n}(x)$ has the following ntegral representation
	\begin{equation}
		\He_{n}(x) = \frac{n!}{2\pi\ii} \int_{|t| = 1} \frac{\ee^{tx - \frac{t^2}{2}}}{t^{n+1}} \dd t
	\end{equation}
	\label{Prop: Hermite Integral Representation}
\end{prop}
\begin{proof}
	We will explicitly arrive at the \textit{Standart} notation for the polynomial using the integral representation
	\begin{equation}
		\He_{n}(x) = \frac{n!}{2\pi\ii} \int_{|t| = 1} \frac{\ee^{tx - \frac{t^2}{2}}}{t^{n+1}} \dd t
	\end{equation}
	Take the Taylor expansion for the exponential
	\begin{align}
		\ee^{tx} & = \sum_{i=0}^{\infty} \frac{(x)^i}{i!} t^i \\
		\ee^{-\frac{t^2}{2}} & = \sum_{j=0}^{\infty} \frac{(-1)^j}{2^j j!} t^{2j}\\
	\end{align}
	and explicitly substitute the expressions into the integral form
	\begin{equation}
		\He_{n}(x) = \frac{n!}{2\pi\ii} \sum_{i=0}^{\infty} \frac{(x)^i}{i!} \sum_{j=0}^{\infty} \frac{(-1)^j}{2^j j!} \int_{|t| = 1} t^{-n-1+i+2j} \dd t
	\end{equation}
	and note that this integral is valued zero except when $2j - (i-n) = 0$. Then, we can write, for every even $i-n$,
	\begin{align}
		\He_n(x) & = n! \sum_{i=0|i-n\text{even}}^{\infty} (-1)^{\frac{n-i}{2}} \frac{x^i}{2^{\frac{n-i}{2}} i! \left( \frac{n-i}{2} \right)!} \\
		& = n! \sum_{l=0}^{\lfloor \frac{n}{2} \rfloor} (-1)^{l} \frac{x^{n-2l}}{2^l l! (n-2l)!}
	\end{align}
	which is exactly the expression given by \ref{Equation: Standart Hermite Notation}.
\end{proof}

\begin{prop}
	I believe there is a connection between the two integral representations if we start with the Gaussian expectation formula, expand the binomial and rewrite the moments with the Cauchy formula, as its done to derive this representation starting from the derivative formula. Missing proof. 
\end{prop}